While the reported results establish a clear solver-level improvement, several limitations should be stated explicitly. First, all endurance evaluations in this paper were conducted on flat terrain in Isaac Sim. This setting is appropriate for isolating the branch-consistency failure mode and for validating the effect of the IPM-based repair under controlled conditions, but it does not exhaust the range of contact phenomena encountered by field robots. In particular, rough terrain and deformable terrain introduce larger contact uncertainty, changing support geometry, and stronger model mismatch, all of which can interact with the optimizer in ways that are not fully exercised by flat-ground experiments. A necessary next step is therefore to extend the proposed IPM + branch-constraint framework to perceptive locomotion over uneven and compliant surfaces, and to evaluate whether the same feasible-set shaping strategy remains robust when contact feedback becomes substantially less predictable.

Second, although the proposed controller preserved stable long-horizon walking, the forward distance did not scale perfectly linearly with time in the 60-second endurance evaluation. The final displacement of 0.638~m confirms persistent forward motion and real-time stability, yet it also indicates that the system gradually adopts a more conservative operating regime over very long horizons. This behavior likely reflects the interaction between heuristic footstep planning, moving-goal reference generation, and small but cumulative state-estimation or target-reconstruction errors. In other words, once the catastrophic branch-switching pathology is removed, a more subtle limitation emerges: the controller remains safe and stable, but its upper-layer motion-generation heuristics may still bias it toward long-horizon stabilization rather than sustained aggressive progression. Addressing this non-linear distance scaling will require improvements in the trajectory generator, foothold policy, and heading-regulation logic, rather than further changes to the branch-consistency constraint itself.

These limitations also help position the broader significance of the paper. Our contribution is strongest at the level of solver correctness and mechanism-aware feasibility: we show that branch validity must be treated as an optimization-domain property rather than as a runtime patch. Future work should build on this corrected solver foundation in two directions. One direction is environmental complexity, including rough, deformable, and partially observed terrain. The other is long-horizon performance shaping, including better footstep heuristics, improved reference generation, and tighter coupling between perception, locomotion intent, and whole-body dynamics. Together, these extensions would test whether the branch-safe IPM formulation introduced here can serve not only as a fix for a flat-ground failure mode, but as a scalable design principle for high-dimensional legged locomotion in realistic settings.
